\section{Discussion}\label{sec:discussion}

\subsection{The applicability boundary as the object of study}

The central empirical fact of this paper is not that any single normalization wins, but that the \emph{sign} of the instance-normalization effect flips across datasets in a way that is predictable before training. The variance attribution of \Cref{sec:empirical} (\Cref{tab:attribution}) makes this precise: the norm main effect explains only $0.73\%$ of log-MSE variance, while the norm$\times$dataset interaction explains $47.12\%$---normalization-related terms total $48.2\%$, against $0.6\%$ for all backbone-related terms. A small main effect with a large interaction is exactly the signature of an applicability boundary: averaged over datasets the normalization choice looks unimportant, yet within any given dataset it is close to decisive. This observation also explains why leaderboard-style comparisons on a fixed benchmark suite can be simultaneously reproducible and uninformative about deployment: whether instance normalization helps depends on a property of the data-generating process, not of the model, echoing recent warnings about benchmark monoculture in long-term forecasting \citep{brigato2026champions,position2025benchmark}.

For practitioners the boundary is actionable because it is cheap to locate. Our recommendation is procedural: before committing to a normalization scheme, compute $\lps$ from the covariates already at hand (\Cref{sec:lps}); if $\lps > \tau$ with $\tau = 0.3$, prefer conditional normalization, otherwise retain instance normalization. The rule is not fragile in the threshold: the pre-registered $8/8$ sign predictions are maintained for every $\tau \in [0.30, 0.70]$ (\Cref{tab:tau}), so the decision requires only that the series fall on the correct side of a wide band, not a finely tuned cut point. We emphasize that this is a characterization, not a rebuttal, of RevIN \citep{kim2022revin}: on the standard benchmarks, where the level is not exogenously predictable, the distribution-shift motivation holds and instance-normalization variants remain in every model confidence set.

\subsection{Capacity and the marginal utility of conditional normalization}

The covariate-fair experiments of Block B refine \emph{what} conditional normalization contributes. When all arms receive identical past and future covariates, the \rawm-versus-\cnm{} gap shrinks monotonically along the mixer capacity ladder: $+0.127$ for the linear mixer, $+0.052$ for the MLP mixer, and $-0.009$ (Raw slightly ahead) for gradient-boosted trees. This is exactly the pattern implied by the reparameterization argument of \citet{toner2024linear}: for a sufficiently unconstrained function class with access to the same information, an input-side normalization is absorbable into the hypothesis space, so its marginal value must vanish. We flag this capacity narrative as an exploratory post-hoc observation, not a pre-registered result, and it is confined to the mixer ladder: PatchTST-Cov ($+0.142$) sits outside the ordering---its minimal covariate adapter exposes less level information to the backbone---and SegRNN-Cov is $-0.004$. The \inm-versus-\cnm{} gap, by contrast, \emph{persists at every capacity}---indeed, RevIN-plus-covariates is worse than \rawm-plus-covariates for every Block~B backbone. This is not a capacity phenomenon but a function-class constraint (\Cref{sec:theory}, Proposition~1): instance normalization projects out the window level before the backbone sees the data, and no amount of downstream flexibility can recover information that the covariates carry about a component that has been subtracted and will be mechanically restored. The refined claim is therefore: the value of \cnm{} is (i) repairing constrained deep pipelines whose built-in normalization discards a predictable level, and (ii) providing a modular route for injecting exogenous level information without redesigning the backbone---not a claim that covariate information is unusable any other way.

\subsection{The cost of conditional normalization when it fails}

The boundary cuts both ways, and the failure mode is not mild. On ETTh1, CondNorm reaches mean MSE $2.2016$ against $0.3771$ for RevIN; on ETTh2, $6.6044$ against $0.2869$ for the best arm (\Cref{tab:main}). The mechanism is first-stage error propagation: when $\lambda \approx 0$ the covariates carry no level signal, the first-stage regression fits noise, and the denormalization step injects that noise directly into the forecast---the $\sest$ term of the \cnm{}-est closed-form risk (\Cref{eq:app-msecn}), which is paid in full precisely when there is no offsetting reduction in the level variance to be explained. The catastrophic magnitudes on the standard benchmarks are thus not an implementation artifact but the predicted behaviour of the estimator outside its applicability region, and they are the reason the decision rule exists at all. A natural safeguard is shrinkage: if the first-stage validation $R^2$ is non-positive, shrink the estimated level $\hat m$ toward the global training mean, which recovers the \rawm{} estimator in the limit. We propose this as an extension and deliberately do \emph{not} apply it post hoc to any number reported in this paper, since doing so would blunt exactly the failure the theory predicts and the pre-registration was designed to test.

\subsection{The persistence blind spot and \texorpdfstring{$\dlps$}{Delta-LPS}}

The absolute $\lps$ has a known blind spot: strongly persistent, covariate-driven series---system marginal prices are the canonical example---can score a high $\lps$ even when the covariates add little beyond what the recent history already implies, because covariates and lagged levels are themselves collinear. The incremental diagnostic $\dlps$ (\Cref{sec:lps}) is the refinement: it measures the covariate contribution \emph{over} a persistence baseline. Our theoretical expectation for this quadrant is regime-dependent: in calm regimes instance normalization should remain adequate (persistence carries the level), while conditional normalization should earn its keep at regime shifts, where the covariates lead and the window statistics lag. We leave this as a targeted stress test rather than asserting it from the present data; the pre-registered rule remains the absolute-$\lps$ threshold, with $\dlps$ reported as a proposed refinement.

\paragraph{Epistemic status of the claims}
It is worth being explicit about which of our statements are theorems,
which are bounds, and which are extrapolations. \emph{Exact}: the
closed-form risks and the crossover identity of
Propositions~1--3 (verified against Monte Carlo to relative error below
$10^{-2}$ and pinned by unit tests), and Proposition~2$'$, which gives the
optimal predictor within each normalization's \emph{linear} function class
as a constrained least-squares solution with no stochastic optimization
involved. \emph{Bound}: the stylized model M1 brackets the crossover from
above; its threshold is an upper bound whose gap to the exact crossing
measures implicit level tracking (\Cref{sec:synthetic}). \emph{Empirical
extrapolation}: no theorem here covers nonlinear backbones. The evidence
that the same geometry governs PatchTST, SegRNN, and the SOTA covariate
architectures is the pre-registered grid and robustness blocks of
\Cref{sec:empirical}, together with a synthetic demonstration that
replacing the linear backbone by a small MLP preserves the crossover
pattern (Fig.~\ref{fig:mlpdemo}, \Cref{app:g7}). Readers should weight the
three tiers accordingly: the boundary's existence and location are proved
for linear classes and measured---not derived---for nonlinear ones.

\subsection{Limitations}

Four limitations bound our claims. First, the theory rests on Assumptions A1--A4 (\Cref{sec:theory}); the synthetic grid shows the constrained-OLS predictions are accurate for linear backbones, but the M1 restoration-rule threshold is an upper bound once backbones implicitly track residual levels. Second, residual design asymmetries remain despite the audit protocol (\Cref{app:design}): the published SegRNN subtracts the last window value rather than the mean, so our revin arm is an approximation of its built-in scheme; PatchTSTCov and SegRNNCov are our own minimal covariate extensions, as no standard recipe exists; the GEFCom2014 load temperature covariate is ex-post practice (identical across arms); and per-block epoch caps differ (Block~A: 12, Block~B: 15), which we control by prohibiting cross-block numeric comparison. Third, the pre-registered results concern deterministic point forecasts under squared error; a post-hoc quantile extension (\Cref{app:g7}) finds the same boundary in pinball loss on both a linear and a gradient-boosted quantile backbone, but also shows that conditional normalization under-covers its nominal intervals when first-stage uncertainty is not propagated---a limitation of the method as instantiated here, flagged for future work. Fourth, our multivariate-SOTA adapters normalize only the target channel, which we argue is conservative toward RevIN because it leaves covariate level information intact; a post-hoc \texttt{revin\_all} arm quantifying the published default (window-normalizing covariate channels as well) supports the choice --- it stays within $\pm 0.05$ MSE of target-only normalization on seven of eight datasets and fails catastrophically on the eighth, where an intermittent precipitation covariate makes window statistics degenerate (\Cref{app:g7}).

\subsection{Implications for power-system operations}

The exogenous group is not incidental: wind, solar, and load are the settings where the boundary matters operationally. Forecast consumers in virtual power plants and system operation care disproportionately about ramps and regime transitions---exactly the events at which a window mean is a stale summary and an archived numerical weather forecast is informative \citep{hong2016gefcom,patil2022wind}. Since $\lps$ requires only the covariates and a first-stage regression already present in most operational pipelines, it can serve as a near-zero-cost screening step: compute it once per series, and let it decide whether the normalization layer should be adaptive to the window or conditional on the forecast drivers.

\section{Conclusion}\label{sec:conclusion}

We have characterized the applicability boundary of instance normalization in multi-step forecasting. Theoretically, we derived closed-form mean-squared errors for the \rawm, \inm, and \cnm{} estimators in a linear-Gaussian model and showed that the benefit of instance normalization decreases monotonically in the exogenous level predictability $\lambda$, with a crossover $\lamstar$ that shrinks as the horizon grows (\Cref{sec:theory}). Empirically, the synthetic grid reproduced the predicted crossings to within $0.015$, and a pre-registered study on eight real datasets---predictions committed before any grid run---achieved $8/8$ sign hits, with the $\lps$--gap association at Spearman $\rho = 0.762$ ($p = 0.028$) at the dataset level and $\rho = 0.783$ per (dataset, horizon) cell. Where the boundary favours conditional normalization it does so decisively: all 11 exogenous model-confidence-set cells retain \{condnorm\} alone, while instance-normalization variants survive on every standard benchmark, where their motivating assumption holds. Practically, the Level Predictability Score converts this boundary into a pre-training decision rule whose conclusions are unchanged for any threshold $\tau \in [0.30, 0.70]$. The broader outlook is a shift from asking \emph{which} normalization is best to diagnosing \emph{when} each is appropriate---we expect similar predictable boundaries for other architectural defaults currently chosen by benchmark averages.
